\chapter*{Resumen}

UMBRAL es un sistema de vigilancia conductual para pequeñas y medianas empresas argentinas que aborda el punto ciego posterior a la autenticación: un atacante con credenciales válidas opera con los permisos del usuario legítimo sin disparar los controles basados en reglas. Siguiendo la separación entre punto de decisión y punto de aplicación de políticas de NIST SP 800-207, un agente liviano para Windows recolecta metadatos de procesos, archivos y conexiones y los envía por TLS a una API central. La API calcula un puntaje de riesgo de 0 a 100 por evento, con continuidad entre sesiones y decaimiento temporal, y devuelve una acción graduada que el agente ejecuta en el equipo: desde una alerta hasta la terminación de procesos o la denegación reversible del acceso a carpetas. Un panel web muestra el riesgo de cada usuario con su explicación y permite ajustar umbrales y acciones en tiempo de ejecución.

El modelo combina Isolation Forest sobre la actividad diaria con una regla explícita de ubicación y una detección de novedad de recurso, entrenados con registros reales del conjunto CLUE-LDS. Sobre días reales, la tasa de días en nivel alto es del 1,81\,\% en el conjunto de desarrollo y del 1,05\,\% en una única evaluación sobre el conjunto reservado, y la latencia de cómputo tiene un percentil 95 de 13,5~ms. Las 795 pruebas automatizadas superan el 85\,\% de cobertura de líneas. La capacidad de detección no pudo medirse, porque ningún conjunto de datos del dominio incluye ataques reales etiquetados, y el modelo entrenado en la nube marca el 53,66\,\% de los días de un dominio de equipos Windows, frente al 4,13\,\% de uno entrenado en ese dominio. Recolectar telemetría propia y obtener incidentes etiquetados son, por eso, el principal trabajo futuro. El análisis económico muestra viabilidad en un escenario de adopción moderada.

\vspace{1em}
\noindent\textbf{Palabras clave:} UEBA, Zero Trust, detección de anomalías, Isolation Forest, puntaje de riesgo, respuesta en el \emph{endpoint}, PyME.
